\section*{Baseline and Evaluation Plan}

The main baseline will be a brute-force collision-detection algorithm that checks every pair of
particles to see if a collision has happened. This method needs $O({n}^{2})$ comparisons for $n$ particles. Will
give a reference point for measuring how much faster spatial partitioning can be. The Quadtree
approach will also be tested against a spatial grid, which splits the simulation area into fixed-size cells
and limits collision checks to nearby cells. All three algorithms will use the collision-detection rules and
particle settings so that the comparison is fair.

Testing will use generated particle placements inside a bounded two-dimensional space. The
placements will cover random positions, clusters, very dense spots and sparse spreads. The number of
particles will grow step by step: 100 first then 500, 1,000, 2,500, 5,000, 10,000, and maybe 25,000 or
more depending on how the hardware handles it. Extra tests will change Quadtree settings, such as how
many particles a node can hold and how deep the tree can grow, to see how the setup changes speed.

The main measurements will be the time taken for collision detection the number of candidate
particle pairs that were checked the time to build the index and the amount of memory used. The
Quadtree will also record how nodes it has and its deepest level so we can see how much overhead the hierarchy
adds. Every algorithm will run with the particle setups and we will repeat each test many times to lower the effect
of random noise. We will confirm correctness by comparing the collision pairs found by each method with those found
by the brute-force baseline.

Results will be shown with tables and graphs that compare how fast each algorithm is across
numbers of particles different spatial layouts and different settings. Especially plots of time and
candidate-pair counts, versus particle count will show how each method scales when the workload
grows. We will also compare how the particle layout and the Quadtree setup affect speed. Finally, we
will analyze the data against what theory predicts for each algorithm so we can spot when spatial
partitioning gives the advantage and when its extra cost makes it less useful.


